\documentclass[12pt,a4paper]{article}
\usepackage[utf8]{vietnam}
\usepackage{amsmath,amssymb,tikz,tkz-tab,multirow,booktabs,xcolor,enumitem}
\usepackage[top=1.5cm,bottom=1.5cm,left=2cm,right=1.5cm]{geometry}
\newcommand{\diem}[1]{\textcolor{blue}{(#1 điểm)}}
\newcommand{\vecto}[1]{\overrightarrow{#1}}
\def\tapR{\mathbb{R}}
\DeclareMathOperator{\Ln}{Ln}
\newtheorem{dl}{Định lí}
\newenvironment{khunglg}{\par\noindent\textbf{Hướng dẫn:}}{\par}
\newif\ifdapan \dapantrue
\title{ĐỀ KIỂM TRA GIỮA KÌ I}
\author{Tổ Toán -- Trường THPT ABC}
\date{}
\begin{document}
\maketitle
\section*{PHẦN I. Trắc nghiệm}
\cau{1} Cho tập hợp $A=\{x\in\tapR \mid x^2-3x+2=0\}$. Số phần tử của $A$ là \diem{0,25}
\bonpa{$2$.}{$3$.}{$5$.}{$4$.}

\cau{2} Giá trị của $\lim\limits_{x\to 0}\dfrac{\sin 3x}{x}$ bằng
\bonpa{$9$.}{$5$.}{$7$.}{$11$.}

\begin{ex}%[0D1Y1-1]
Trong mặt phẳng $Oxy$, cho $\vecto{AB}=(1;2)$ và $\vv{AC}=(3;-1)$. Tính $\left\lvert\vecto{AB}\cdot\vv{AC}\right\rvert$.
\choice{$1$}{\True $2$}{$3$}{$4$}
\loigiai{Ta có $\vecto{AB}\cdot\vv{AC}=1\cdot 3+2\cdot(-1)=1$, suy ra \textbf{đáp án} là $1$.}
\end{ex}

\begin{ex}
Mệnh đề nào sau đây \textbf{đúng}?
\choiceTF{\True Hàm số $y=x^3$ đồng biến trên $\mathbb{R}$}{Phương trình $x^2+1=0$ có nghiệm}{\True $\sqrt{4}=2$}{$\pi$ là số hữu tỉ}
\end{ex}

\section{Tự luận}
\begin{dl}[Pytago]
Tam giác $ABC$ vuông tại $A$ khi và chỉ khi $BC^2=AB^2+AC^2$.
\end{dl}

\begin{enumerate}[label=\alph*)]
\item Giải phương trình $\heva{&x+y=3\\&x-y=1}$.
\item Tính tích phân
\begin{align}
I &= \int_0^1 x e^{x}\,\mathrm{d}x \label{eq:I}\\
  &= \Big[x e^x\Big]_0^1 - \int_0^1 e^x\,dx = 1. \nonumber
\end{align}
Theo \eqref{eq:I} ta có kết quả.
\item Chứng minh $\Ln(xy)=\Ln x+\Ln y$ với $x,y\in\mathbb{R}^{+}$ và $a\not\in B$, $\overparen{AB}$.
\end{enumerate}

\begin{itemize}
\item Văn bản có chú thích\footnote{Đây là chú thích chân trang với $x^2$.} và \emph{nhấn mạnh}.
\item Màu \textcolor{red}{chữ đỏ} và \colorbox{yellow}{nền vàng}, \underline{gạch chân}, \sout{gạch ngang}.
\end{itemize}

\begin{center}
\begin{tabular}{|c|c|c|}
\hline
\multirow{2}{*}{Nhóm} & \multicolumn{2}{c|}{Kết quả} \\
\cline{2-3}
 & Đạt & Không đạt \\
\hline
A & $12$ & $3$ \\
\hline
\end{tabular}
\end{center}

\begin{table}[h]
\centering
\caption{Bảng số liệu}\label{tab:1}
\begin{tabular}{lcr}
\toprule
\textbf{Tên} & Tuổi & Điểm \\
\midrule
An & 16 & 8,5 \\
Bình & 17 & 9,0 \\
\bottomrule
\end{tabular}
\end{table}
Xem Bảng~\ref{tab:1}.

\ifdapan
\begin{khunglg}
Đây là phần hướng dẫn chỉ hiện khi bật đáp án.
\end{khunglg}
\fi

Họ và tên: \dotfill\ Lớp: \hrulefill

\begin{verbatim}
\begin{document} 100% $x$
\end{verbatim}

Hoá học: \ce{2H2 + O2 -> 2H2O}, đơn vị \SI{9.8}{m/s^2}.
\end{document}
